\documentclass[12pt]{article}
\usepackage{tikz}
\usepackage{float}
\usepackage{amsmath, amssymb}
\usepackage{amsmath, amssymb}
\usetikzlibrary{shapes, backgrounds, calc, positioning, decorations.pathreplacing, shadows, arrows.meta}

% --- Colors ---
\definecolor{realColor}{RGB}{0, 100, 200}    % Deep Blue
\definecolor{synColor}{RGB}{220, 50, 50}     % Deep Red
\definecolor{transportColor}{RGB}{230, 140, 0} % Orange for Wasserstein
\usetikzlibrary{shapes.geometric, arrows.meta, positioning, calc, backgrounds, fit, shadows}
\usepackage{amsmath, amssymb}
\usetikzlibrary{shapes, backgrounds, calc, positioning, decorations.pathreplacing, shadows}

% --- Colors ---
\definecolor{realColor}{RGB}{0, 100, 200}    % Deep Blue
\definecolor{synColor}{RGB}{220, 50, 50}     % Deep Red
\definecolor{overlapColor}{RGB}{100, 0, 100} % Purple (Overlap)

% --- Professional Colors ---
\definecolor{encoderColor}{RGB}{230, 240, 250} % Light Blue
\definecolor{decoderColor}{RGB}{250, 230, 230} % Light Red
\definecolor{tensorColor}{RGB}{255, 255, 255}  % White
\definecolor{lineColor}{RGB}{50, 50, 50}       % Dark Grey
\definecolor{highlight}{RGB}{0, 100, 180}      % Deep Blue
% ---- Select text fonts ----
\input{glyphtounicode}
\pdfgentounicode=1
\usepackage[T1]{fontenc}
\usepackage{sourceserifpro}
\usepackage{sourcecodepro}
\usepackage[hyphens]{url}

% ---- Select math fonts ----
\usepackage{MnSymbol}
\usepackage[cal=euler,bb=fourier]{mathalpha}
\newcommand{\hmmax}{0}
\newcommand{\bmmax}{0}
\usepackage{bm}
\usepackage[italic,eulergreek,nolessnomore,noparenthesis,noplusnominus,noequal,nohbar]{mathastext}
\MTsetmathskips{f}{3mu}{0mu}
\MTsetmathskips{y}{2mu}{0mu}
\MTsetmathskips{p}{2mu}{0mu}
\MTsetmathskips{j}{2mu}{0mu}

% ---- Page layout ----
\usepackage{geometry}
\geometry{top=1in, bottom=1.2in, left=1.3in, right=1.3in, footskip=0.7in}

% ---- Spacing ----
\usepackage[onehalfspacing]{setspace}
\usepackage{microtype}
\AtBeginDocument{\frenchspacing}

% ---- Colors ----
\usepackage[x11names]{xcolor}

% ---- Headers and footers ----
\usepackage{fancyhdr}
\pagestyle{fancy}
\renewcommand{\headrule}{}
\fancyhf{}
\fancypagestyle{plain}{\fancyhf{}}
\fancyfoot[C]{\thepage}

% ---- Title page ----
\usepackage{titling}
\usepackage{xparse}
\pretitle{\begin{center}\bfseries\Huge}
\renewcommand{\tamark}{}
\setlength{\thanksmarkwidth}{0em}
\setlength{\thanksmargin}{0em}

\NewDocumentCommand{\available}{o g}{%
\fancypagestyle{plain}{\fancyfoot[C]{\color{gray}\footnotesize%
\IfValueT{#1}{Published in \emph{#1}\;\textbullet\;}%
Available at \url{#2}}}}

\usepackage{etoolbox}
\AtBeginEnvironment{titlepage}{\pagestyle{empty}\setlength{\parindent}{0pt}}

% ---- Section headings ----
\usepackage{titlesec}
\titleformat*{\section}{\centering\large\bfseries}
\titleformat*{\subsection}{\bfseries}
\titleformat{\paragraph}[runin]{\itshape}{}{}{}[\@addpunct{.}]
\titlelabel{\thetitle.\quad}

% ---- Theorems and proofs ----
\usepackage{amsthm}
\newtheoremstyle{paper}{}{}{\itshape}{}{\scshape}{.}{0.5em}{}
\theoremstyle{paper}
\newtheorem{theorem}{theorem}
\newtheorem{proposition}{proposition}
\newtheorem{lemma}{lemma}
\newtheorem{corollary}{corollary}
\newtheorem{definition}{definition}
\newtheorem{assumption}{assumption}
\newtheorem{remark}{remark}
\renewcommand{\proofname}{\upshape\scshape proof}

% ---- Figures and tables ----
\usepackage{graphicx}
\usepackage{booktabs}
\usepackage{rotating}
\usepackage{multirow}
\renewcommand{\floatpagefraction}{0}
\makeatletter\g@addto@macro\@floatboxreset\centering\makeatother
\renewcommand{\arraystretch}{1.1}
\BeforeBeginEnvironment{tabular*}{\footnotesize}
\renewcommand{\tablename}{\textsc{table}}
\renewcommand{\figurename}{\textsc{figure}}

\usepackage{caption,subcaption}
\captionsetup{labelsep=period}
\captionsetup[table]{position=top}
\captionsetup[sub]{labelformat=simple,labelsep=period,size=small}
\renewcommand{\thesubfigure}{\Alph{subfigure}}

\NewDocumentCommand{\note}{o g}{%
\parbox{\textwidth}{\footnotesize\vspace*{10pt}%
\IfValueT{#1}{\textit{#1}:\quad}%
#2}}

% ---- Equations ----
\usepackage{amsmath}
\AtBeginDocument{\allowdisplaybreaks[1]}

% ---- Lists ----
\usepackage{enumitem}
\setlist[itemize,1]{leftmargin=\parindent,label=\color{gray}{\upshape\textbullet}}
\setlist[itemize,2]{leftmargin=2\parindent,label=\color{gray}{\upshape\textendash}}
\setlist[enumerate,1]{leftmargin=\parindent,label=\color{gray}{\upshape\alph*.}}
\setlist[enumerate,2]{leftmargin=2\parindent,label=\color{gray}{\upshape\roman*.}}
\setlist{nosep}

% ---- Bibliography ----
\usepackage{natbib}
\setcitestyle{aysep={}}
\renewcommand{\bibfont}{\small}
\setlength{\bibsep}{0pt}
\setlength{\bibhang}{\parindent}

% ---- Appendix formatting ----
\AfterEndEnvironment{thebibliography}{
\setcounter{theorem}{0}
\setcounter{proposition}{0}
\setcounter{lemma}{0}
\setcounter{corollary}{0}
\setcounter{definition}{0}
\setcounter{assumption}{0}
\setcounter{remark}{0}
\setcounter{table}{0}
\setcounter{figure}{0}
\setcounter{equation}{0}
\renewcommand{\thetheorem}{A\arabic{theorem}}
\renewcommand{\theproposition}{A\arabic{proposition}}
\renewcommand{\thelemma}{A\arabic{lemma}}
\renewcommand{\thecorollary}{A\arabic{corollary}}
\renewcommand{\thedefinition}{A\arabic{definition}}
\renewcommand{\theassumption}{A\arabic{assumption}}
\renewcommand{\theremark}{A\arabic{remark}}
\renewcommand{\thetable}{A\arabic{table}}
\renewcommand{\thefigure}{A\arabic{figure}}
\renewcommand{\theequation}{A\arabic{equation}}
\titleformat{\section}{\centering\large\bfseries}{Appendix~\thesection.}{1em}{}}

% ---- Hyperlinks ----
\usepackage[pdftex,hidelinks,hypertexnames=false,pdfpagemode=UseNone,pdfdisplaydoctitle=true]{hyperref}

% =========================================================================
% DOCUMENT METADATA
% =========================================================================

\title{Synthetic Data Augmentation in Glioma Raman Spectroscopy}
\author{Andrei Iușan, Iulian Vasile}
\date{\today}

% Optional: Add publication info to footer
% \available[Journal Name]{https://yourdomain.com/paper}

% =========================================================================
% DOCUMENT CONTENT
% =========================================================================

\begin{document}

\maketitle

\begin{abstract}
Access to large-scale, statistically representative biomedical datasets is critical for advancing AI-driven diagnostics, particularly for heterogeneous pathologies like glioma. However, the acquisition of Raman spectroscopy data is often hindered by high costs, time constraints, and strict patient privacy regulations, resulting in small sample sizes that lead to model overfitting. To address these challenges, we introduce a 1D-Convolutional Conditional Variational Autoencoder (CVAE) for generating high-fidelity synthetic Raman spectra. Unlike traditional oversampling methods, our deep learning approach learns the underlying non-linear statistical manifold of biochemical features, allowing for the generation of infinite, diverse spectral signatures conditioned on tumor subtype. We validate the utility of this approach using a 'Train on Synthetic, Test on Real' (TSTR) framework. Our results demonstrate that synthetic augmentation not only preserves the biological fidelity of the original data but also significantly improves the robustness of downstream classifiers in data-scarce environments
\end{abstract}

\section{Introduction}
\subsection{Clinical Context}
\label{sec:clinical_context}
The current gold standard for intraoperative glioma grading relies on frozen section histopathology. While effective, this process is time-consuming (days for lab test), requires tissue excision and wait time for results, and is subject to inter-observer variability. In contrast, Raman Spectroscopy (RS) offers a non-destructive, label-free optical biopsy capable of detecting biochemical fingerprints of tissue in seconds. By targeting specific molecular vibrations---such as the phenylalanine peak at $1003 \text{ cm}^{-1}$ or lipid-protein ratios at $1450 \text{ cm}^{-1}$---RS can theoretically delineate tumor margins with high precision. This would allow for real-time intraoperative decision-making. However, the translation of RS into clinical workflows is bottlenecked by the difficulty of training robust machine learning classifiers due to high signal noise, batch effects, and limited sample sizes.

\subsection{Challenges in Raman Spectral Analysis}
\label{sec:challenges}
Despite its clinical potential, the adoption of AI-assisted Raman spectroscopy into surgical workflows is bottlenecked by the inherent complexities of spectral data. Biomedical spectral datasets typically suffer from the ``curse of dimensionality,'' operating in a ``small $N$, large $P$'' regime[cite: 14]. A typical patient cohort might consist of $N < 100$ subjects, yet each Raman spectrum contains $P > 1000$ spectral bands[cite: 15]. When standard machine learning classifiers (e.g., Support Vector Machines, Random Forests) are trained on such high-dimensional, low-sample data, they are prone to severe overfitting[cite: 16]. Rather than learning generalizable biochemical features, these models tend to memorize patient-specific noise[cite: 16].

This overfitting is exacerbated by a severe domain shift problem. Raw Raman spectra are frequently contaminated by fluorescence background, cosmic ray artifacts, and instrumental noise that vary between acquisition sessions and spectrometers[cite: 33]. In heterogeneous pathologies like glioma, the subtle biochemical differences between sub-grades are often drowned out by those hardware and aquisition factors; the session-dependent variance ($\sigma_{\text{batch}}^2$) frequently exceeds the biologically meaningful variance ($\sigma_{\text{biology}}^2$)[cite: 34]. 

Crucially, the magnitude of this generalization gap is often masked in the current literature due to optimistic bias. Even when strict patient-level cross-validation is enforced, the ``small $N$, large $P$'' regime makes machine learning pipelines susceptible to hyperparameter overfitting. Extensive tuning of feature selection algorithms, clustering parameters, and model architectures on a fixed dataset can artificially inflate diagnostic accuracies, often pushing binary classification performance just past the 80\% threshold. While this does not constitute data leakage, this hyperparameter optimization creates brittle models and overestimates the true generalization performance on unseen data. Those models also fail as the clinical complexity increases from binary classification to multi-class grading of glioma subtypes, where the spectral differences are more subtle and noise is more likely to dominate the signal.

\subsection{Proposed Solution and Contributions}
\label{sec:contributions}
This study aims to bridge the generalization gap by developing a robust generative framework that disentangles underlying tumor biology from instrument-induced batch effects. We hypothesize that deep generative models can not only synthesize high-fidelity training data to regularize downstream classifiers in data-scarce environments, but that their generative physics can be directly leveraged for zero-shot inference. Our specific contributions are threefold:

\begin{enumerate}
    \item \textbf{Spatio-Spectral Tumor Annotation:} We introduce an unsupervised, per-slide annotation pipeline that combines spectral dimensionality reduction (PCA) with spatial coordinate weighting. By fitting a Gaussian Mixture Model (GMM) to individual biopsy slides, we successfully isolate pure tumor spectra from healthy tissue while entirely bypassing the confounding influence of global batch-effect clustering.
    \item \textbf{Deep Generative Augmentation ($\beta$-CVAE):} We design a 1D-Convolutional Conditional Variational Autoencoder optimized to learn the non-linear statistical manifold of Raman spectra. By conditioning the latent space on six distinct glioma subtypes, we generate infinite, perfectly balanced synthetic datasets. We validate the utility of this approach by demonstrating massive performance gains in a Train-on-Synthetic, Test-on-Real (TSTR) framework compared to strictly isolated Train-on-Real, Test-on-Real (TRTR) baselines.
    \item \textbf{Classification by Reconstruction (CbR):} We propose a novel generative inference paradigm that completely bypasses standard discriminative decision boundaries. By encoding a raw test spectrum into its pure physical representation ($z$) and decoding it across all possible class conditions, we classify the tumor based on the lowest Mean Squared Error (MSE) of the reconstruction---matching the performance of heavily optimized supervised baselines using pure physics.
\end{enumerate}



\section{Background and Related Work}
\label{sec:background}

\subsection{Deep Generative Augmentation ($\beta$-CVAE)}
\label{sec:cvae_architecture}
To learn the underlying statistical manifold of the pure tumor spectra, we implemented a 1D-Convolutional Conditional Variational Autoencoder ($\beta$-CVAE). Rather than relying on generic multi-layer perceptrons, 1D convolutions were chosen to explicitly exploit the strong local collinearity between adjacent Raman shift frequencies (e.g., the broad lipid/protein bands between $1400$ and $1500 \text{ cm}^{-1}$). 

As illustrated in Figure \ref{fig:architecture}, the architecture is composed of an Encoder $q_\phi(z|x)$, a 32-dimensional latent bottleneck, and a Decoder $p_\theta(x|z,y)$. The encoder compresses the 1738-dimensional input into a latent representation $z$. To ensure the latent space is organized by tumor biology rather than arbitrary hardware noise, we condition the decoder on a one-hot encoded class label $y$, representing the six glioma sub-grades. 

We frame the encoding process as a Multi-Task Learning (MTL) problem by appending an auxiliary classification head to the encoder. The model is optimized using a composite loss function:
\begin{equation}
\mathcal{L}_{\text{total}} = \mathcal{L}_{\text{recon}} + \beta \cdot D_{KL}(q_\phi(z|x) \| p(z)) + \alpha \cdot \mathcal{L}_{\text{class}}
\end{equation}
where $\mathcal{L}_{\text{recon}}$ is the Mean Squared Error (MSE) of the spectral reconstruction. Following the $\beta$-VAE framework introduced by Higgins et al. (2017), the Kullback-Leibler (KL) divergence acts as a structural regularizer scaled by $\beta=2$. This heavily constrains the latent capacity, forcing the network to disentangle the robust biological signal from the high-variance hardware noise. Simultaneously, $\mathcal{L}_{\text{class}}$ is the categorical cross-entropy of the auxiliary classifier scaled by $\alpha = 0.2$. This joint objective prevents the convolutional filters from overfitting to patient-specific noise, forcing the network to extract features that are both physically descriptive enough for reconstruction and biologically distinct enough for classification. Once trained, the decoder can generate infinite, perfectly balanced synthetic datasets by sampling $z \sim \mathcal{N}(0, I)$ and conditioning on a requested class label.

\subsection{Generative Inference: Classification by Reconstruction (CbR)}
\label{sec:cbr}
Standard discriminative models (e.g., Random Forests, XGBoost) attempt to learn rigid mathematical boundaries between classes. As demonstrated by the domain shift problem, these boundaries often fail when applied to unseen patients with varying spectrometer baselines. To bypass this limitation, we propose a novel, zero-shot generative inference paradigm: \textit{Classification by Reconstruction} (CbR).

CbR leverages the CVAE's learned physics to classify a tumor without relying on discriminative boundaries. During inference on a raw, unseen test spectrum $x_{\text{test}}$, the process is defined as follows:
\begin{enumerate}
    \item \textbf{Latent Extraction:} The test spectrum is passed through the Encoder to extract its pure physical shape vector, $\mu_{z} = \text{Encoder}(x_{\text{test}})$.
    \item \textbf{Conditional Hallucination:} The physical vector $\mu_{z}$ is passed to the Decoder $N_c$ times (where $N_c = 6$ classes). Each forward pass pairs $\mu_{z}$ with a different one-hot class condition $y_i$, generating a set of hallucinated spectra: $\hat{x}_i = \text{Decoder}(\mu_{z}, y_i)$.
    \item \textbf{Error Minimization:} We compute the Mean Squared Error between the raw test spectrum and each of the six hallucinations. The final predicted class $\hat{y}$ is the condition that yields the lowest reconstruction error:
    \begin{equation}
    \hat{y} = \underset{i \in \{0, \dots, 5\}}{\mathrm{argmin}} \left( \frac{1}{P} \sum_{j=1}^{P} (x_{\text{test}, j} - \hat{x}_{i, j})^2 \right)
    \end{equation}
\end{enumerate}
This approach posits that the Decoder can only perfectly reconstruct a spectrum if the requested biological class condition matches the true physics encoded in $\mu_{z}$. If forced to reconstruct a Grade IV spectrum using a Grade II condition, the model will hallucinate incorrect biochemical peaks, resulting in a high MSE. CbR thus reframes multi-class grading as an unsupervised signal-matching problem.

\subsection{Variational Autoencoders (VAEs)}
Variational Autoencoders (VAEs) operate as a generative framework that bridges deep learning with Bayesian inference. The core objective is to approximate the true data distribution $p_{\text{data}}(x)$ by utilizing a latent-variable model. This involves a learned decoder, $p_{\theta}(x|z)$, which maps a compressed latent representation $z$---typically governed by a standard Gaussian prior $p(z)$---back into the original data space.

Because calculating the true posterior $p(z|x)$ is computationally prohibitive, VAEs utilize an inference network (the encoder), denoted as $q_{\phi}(z|x)$, to approximate it. The model is optimized by maximizing the \textbf{Evidence Lower Bound (ELBO)}, which serves as a surrogate for the marginal log-likelihood:

\begin{equation}
\mathcal{L}(\theta, \phi; x) = \mathbb{E}_{q_{\phi}(z|x)}[\log p_{\theta}(x|z)] - \beta D_{KL}(q_{\phi}(z|x) \parallel p(z))
\end{equation}

\subsection*{Deconstructing the Objective}

The ELBO objective is comprised of two distinct functional components:
\begin{itemize}
    \item \textbf{Reconstruction Accuracy:} The expected log-likelihood term $\mathbb{E}_{q_{\phi}(z|x)}[\log p_{\theta}(x|z)]$ ensures that the decoder can faithfully reconstruct the input from the latent code.
    \item \textbf{Latent Regularization:} The Kullback-Leibler (KL) divergence term acts as a constraint, penalizing the approximate posterior if it drifts too far from the prior distribution $p(z)$.
\end{itemize}


\subsection{Conditional VAEs}
Standard VAEs generate data from a latent space $z \sim \mathcal{N}(0, I)$, but they lack control over specific semantic attributes of the output. For clinical applications, generating a random spectrum is insufficient; we require the ability to request a specific pathology (e.g., ``Generate a Grade IV Glioblastoma spectrum'').

To achieve this, we employ a Conditional VAE (CVAE). The CVAE modifies the probabilistic framework by conditioning both the encoder and decoder on a class label $y$ (one-hot encoded vector representing tissue type). The objective function becomes a conditional ELBO:
\begin{equation}
\mathcal{L}_{\text{CVAE}}(\theta, \phi; x, y) = \mathbb{E}_{q_{\phi}(z|x, y)}[\log p_{\theta}(x|z, y)] - \beta D_{KL}(q_{\phi}(z|x, y) \parallel p(z|y))
\end{equation}
 $y$ is concatenated with the latent vector $z$ during decoding, guiding the generation process to a specific region of the data manifold associated with the desired pathology.

\subsection{Synthetic Data for Biomedical Applications}
\label{sec:synth_background}
Synthetic data generation has emerged as a promising strategy for overcoming data scarcity in clinical machine learning \cite{jordon2022synthetic}. Unlike classical augmentation techniques (e.g., additive noise, SMOTE), deep generative models learn the underlying data manifold and can produce novel, diverse samples that respect the statistical structure of the original distribution.

The standard evaluation framework for synthetic data utility is the \textbf{Train on Synthetic, Test on Real (TSTR)} paradigm \cite{esteban2017real}. In TSTR, a downstream classifier is trained \emph{exclusively} on generated data and evaluated on held-out real data. If the TSTR accuracy approaches the \textbf{Train on Real, Test on Real (TRTR)} baseline, it demonstrates that the generative model has captured the discriminative features of the real distribution without memorizing individual training examples.

Beyond utility, synthetic data evaluation requires metrics for \textbf{fidelity} (do synthetic samples look realistic?), \textbf{diversity} (does the generator cover the full data manifold or suffer from mode collapse?), and \textbf{privacy} (are synthetic samples sufficiently distant from real training examples?) \cite{alaa2022faithful}.

\section{Methods}
\label{sec:methods}

\subsection{Data and Preprocessing}
\label{sec:preprocessing}
The dataset consists of Raman spectra acquired from fresh glioma biopsies. To ensure the model learns chemical features rather than instrumental artifacts, we applied a standardized preprocessing pipeline:
\begin{itemize}
    \item \textbf{Cosmic Ray Removal:} Outliers exceeding $5\sigma$ of the local neighborhood were excised.
    \item \textbf{Baseline Correction:} Fluorescence background was estimated and removed using Adaptive Iteratively Reweighted Penalized Least Squares (airPLS) with $\lambda=10^4$.
    
    \item \textbf{Vector Normalization:} Each spectrum was normalized using Min-Max scaling to the $[0,1]$ range to standardize intensity variations caused by laser focus depth.
\end{itemize}
Data splitting was performed at the \textit{patient level} using a fixed 5-fold cross-validation scheme, ensuring that spectra from the same tissue sample never appear in both training and testing sets. The splits were frozen to JSON files to guarantee reproducibility across all experiments. In addition, a subset of 6 patients (approximately 10\% of the dataset) was held out as a final test set for zero-shot inference experiments.

\subsection{Spatio-Spectral Tumor Annotation}
\label{sec:annotation}
A critical preprocessing step is separating tumor tissue from healthy tissue within each biopsy scan. Raw hyperspectral cubes contain a mixture of both tissue types, and training generative models on mixed data would dilute the learned class-specific distributions.

We introduce an unsupervised \textbf{Spatio-Spectral Annotation} pipeline that operates on a per-slide basis to bypass global batch effects:
\begin{enumerate}
    \item \textbf{Spectral Dimensionality Reduction:} PCA is applied to the preprocessed spectra of each slide independently, retaining 99\% of the variance.
    \item \textbf{Spatial Context Injection:} The $(x, y)$ pixel coordinates of each spectrum are scaled and concatenated with the PCA features, weighted by a spatial factor $\omega = 2.0$.
    \item \textbf{Gaussian Mixture Model (GMM):} A 2-component GMM is fitted to the combined spatio-spectral features. The cluster with the higher mean phenylalanine signal ($1003$ cm$^{-1}$) is identified as tumor tissue.
\end{enumerate}

This per-slide approach avoids the pitfall of global clustering, where batch effects between slides would dominate the spectral variance and obscure the true biological signal \cite{reynolds2009gaussian}. The resulting binary tumor labeling yields approximately 88,000 pure tumor spectra from the original $\sim$270,000 total.

\subsection{Model Architecture}
We implemented a 1D-Convolutional CVAE to exploit local correlations between adjacent spectral bands. 
\begin{itemize}
    \item \textbf{Encoder:} Consists of three Conv1D layers (filters: 32, 64, 128; kernel size: 3) with LeakyReLU activation and Batch Normalization. This compresses the 1167-dimensional input into a 16-dimensional latent space.
    \item \textbf{Latent Space:} We utilized the reparameterization trick $z = \mu + \sigma \odot \epsilon$ to allow backpropagation.
    \item \textbf{Decoder:} Mirrors the encoder using Conv1D-Transpose layers to reconstruct the original spectral dimensions.
\end{itemize}
The architecture and data flow are visualized in Figure \ref{fig:architecture}.



\begin{tikzpicture}[
    scale=1.05,
    transform shape,
    node distance=1.1cm,
    >={Latex[width=1.5mm,length=1.5mm]},
    font=\sffamily\footnotesize,
    thick,
    tensor/.style={
        draw=lineColor,
        fill=tensorColor,
        rectangle,
        rounded corners=1.5pt,
        minimum height=1.5cm,
        minimum width=0.45cm
    },
    shorttensor/.style={
        draw=lineColor,
        fill=tensorColor,
        rectangle,
        rounded corners=1.5pt,
        minimum height=1.0cm,
        minimum width=0.45cm
    },
    block/.style={
        draw=lineColor,
        rectangle,
        rounded corners=3pt,
        minimum width=2.1cm,
        minimum height=2.8cm,
        align=center
    },
    op/.style={
        draw=lineColor,
        circle,
        fill=white,
        inner sep=0pt,
        minimum size=0.45cm
    },
    head/.style={
        draw=lineColor,
        rectangle,
        fill=yellow!10,
        rounded corners=3pt,
        minimum width=2.0cm,
        minimum height=1.0cm,
        align=center
    }
]

% ==========================================================
% 1. INPUT
% ==========================================================
\node[tensor, label=above:Input $x$] (x) at (0,0) {};
\node[below=0.01cm of x, font=\tiny, color=gray] {$(1167,1)$};

% ==========================================================
% 2. ENCODER
% ==========================================================
\node[block, right=1.0cm of x, fill=encoderColor,
      label={[yshift=-3.5cm]\textbf{Encoder} $q_\phi(z|x)$}] (encoder) {
    \textbf{Conv1D}\\
    $\downarrow$\\
    Flatten\\
    $\downarrow$\\
    Dense
};

% ==========================================================
% 3. LATENT SPACE
% ==========================================================
\node[shorttensor, right=1.1cm of encoder, yshift=0.4cm, label=above:$\mu$] (mu) {};
\node[shorttensor, right=1.1cm of encoder, yshift=-0.8cm, label=below:$\log \sigma^2$] (sigma) {};

\node[op, right=0.9cm of sigma] (mult) {$\otimes$};
\node[below=0.5cm of mult, align=center, font=\scriptsize] (noiselabel)
    {$\epsilon \sim \mathcal{N}(0,I)$};
\draw[->, dashed] (noiselabel) -- (mult);

\node[op] (add) at (mu -| mult) {$\oplus$};
\node[shorttensor, right=0.6cm of add, fill=green!10, label=above:Latent $z$] (z) {};
\node[below=0.1cm of z, font=\tiny, color=gray] {$(16)$};

% ==========================================================
% 4. CLASSIFIER HEAD (branching upward from encoder)
% ==========================================================
\node[head, above=1.2cm of encoder] (classifier) {Classifier\\Head};
\node[tensor, right=1.0cm of classifier, fill=yellow!10, label=above:$\hat{y}$ (Pred)] (ypred) {};

% ==========================================================
% 5. LABEL y (below latent z, feeds into decoder)
% ==========================================================
\node[shorttensor, below=1.2cm of z, fill=red!10, label=below:Label $y$] (y) {};
\node[below=0.45cm of y, font=\tiny, color=gray] {$(N_c)$};

% ==========================================================
% 6. DECODER
% ==========================================================
\node[op, right=0.7cm of z, yshift=-0.6cm] (concat2) {$\oplus$};

\node[block, right=0.6cm of concat2, fill=decoderColor,
      label={[yshift=-3.5cm]\textbf{Decoder} $p_\theta(x|z,y)$}] (decoder) {
    Dense\\
    $\downarrow$\\
    Reshape\\
    $\downarrow$\\
    \textbf{Conv1D-T}
};

% ==========================================================
% 7. OUTPUT
% ==========================================================
\node[tensor, right=1.0cm of decoder, label=above:$\hat{x}$] (xhat) {};
\node[below=0.1cm of xhat, font=\tiny, color=gray] {$(1167,1)$};

% ==========================================================
% 8. CONNECTIONS
% ==========================================================

% Input -> Encoder
\draw[->] (x) -- (encoder);

% Encoder -> mu, sigma
\draw[->] (encoder.east) -- (mu.west);
\draw[->] (encoder.east) -- (sigma.west);

% Reparameterisation
\draw[->] (sigma) -- (mult);
\draw[->] (mult) -- (add);
\draw[->] (mu)   -- (add);
\draw[->] (add)  -- (z);

% Encoder -> Classifier Head (branch upward)
\draw[->] (encoder.north) -- (classifier.south);
\draw[->] (classifier)    -- (ypred);

% z -> concat -> Decoder
\draw[->] (z.east) -| (concat2.north);

% y (label) -> concat (red, annotated)
\draw[->, color=red] (y.east) -- ++(0.3,0)
    node[midway, above, font=\scriptsize, color=red] {Concatenate}
    |- (concat2.south);

% Decoder -> output
\draw[->] (concat2) -- (decoder);
\draw[->] (decoder) -- (xhat);

\end{tikzpicture}
\subsection{Training Protocol}
\label{sec:training}
The CVAE is trained using a composite loss function:
\begin{equation}
\mathcal{L}_{\text{total}} = \mathcal{L}_{\text{recon}} + \beta \cdot D_{KL}(q_\phi(z|x,y) \| p(z)) + \alpha \cdot \mathcal{L}_{\text{class}}
\end{equation}


\subsection{Experimental Design}
\label{sec:experimental_design}
Our evaluation follows a three-tier comparison framework, with all experiments conducted within the same 5-fold patient-level cross-validation:

\paragraph{Baseline (TRTR)} We reproduce the ``Apollo'' methodology from prior work \cite{livanos2024glioma}: a Random Forest classifier (300 trees, entropy criterion) selects the \textbf{Top 20} most discriminative spectral bands, followed by an SVM with RBF kernel ($C=1.0$, $\gamma=\text{scale}$, balanced class weights) trained on these features. This establishes the real-data accuracy ceiling. We also tested RF, XGBoost and a 1D-CNN classifier.

\paragraph{TSTR Evaluation} For each cross-validation fold, the CVAE is trained on the training split. We then generate $N_{\text{syn}}$ synthetic spectra per class from the trained decoder by sampling $z \sim \mathcal{N}(0, I)$ and conditioning on each class label. The \emph{same} suite of classifiers is trained exclusively on synthetic data and tested on the held-out real fold.


\section{Results}
\label{sec:results}

\subsection{Baseline Performance (TRTR)}
\label{sec:trtr_results}
The Apollo pipeline (RF Top-20 + SVM) was evaluated using the frozen 5-fold patient-level cross-validation. Results are presented in Table~\ref{tab:trtr}.

\begin{table}[H]
\caption{TRTR Baseline: Apollo Pipeline (6-class glioma grading)}
\label{tab:trtr}
\begin{tabular*}{\textwidth}{@{\extracolsep\fill}lcc}
\toprule
Fold & Held-Out Patients & Accuracy \\
\midrule
Fold 0 & 14 samples & 0.373 \\
Fold 1 & 12 samples & 0.263 \\
Fold 2 & 14 samples & 0.319 \\
Fold 3 & 10 samples & 0.288 \\
Fold 4 & 12 samples & 0.233 \\
\midrule
\textbf{Mean $\pm$ Std} & & $\mathbf{0.295 \pm 0.051}$ \\
\bottomrule
\end{tabular*}
\note{All folds use strict patient-level isolation. No spectra from the same tissue sample appear in both training and testing sets.}
\end{table}

The mean accuracy of $29.5\% \pm 5.1\%$ across the 6-class task confirms the severity of the generalization gap. While substantially above the $16.7\%$ chance level, this result demonstrates that batch effects dominate inter-class variance when patients are strictly isolated, consistent with the hypothesis outlined in Section~\ref{sec:data_scarcity}.

\subsection{Statistical Fidelity of Synthetic Spectra}
\label{sec:fidelity}

% TODO: Fill with actual Wasserstein distances and KS test results after running the evaluation notebook
To assess the fidelity of generated spectra, we compute the Wasserstein-1 distance between real and synthetic distributions projected onto the 64-dimensional feature space of the 1D-CNN encoder. Lower values indicate that the synthetic manifold better approximates the real data support.

\begin{figure}[H]
\begin{tikzpicture}[font=\sffamily, thick]
    
    % ==========================================================
    % 1. REAL MANIFOLD (The Support)
    % ==========================================================
    \filldraw[
        draw=realColor, 
        fill=realColor, 
        fill opacity=0.1, 
        line width=1.5pt
    ] plot[smooth cycle, tension=0.7] coordinates {
        (-3, 0) (-2, 2.5) (1, 3) (3, 1) (2, -2) (-1, -3) (-3, -1.5)
    };
    
    % Label for Real
    \node[text=realColor, font=\bfseries] at (-4.5, 2.5) {Real Support ($\mathcal{S}_r$)};
    
    % ==========================================================
    % 2. SYNTHETIC MANIFOLD
    % ==========================================================
    \filldraw[
        draw=synColor, 
        fill=synColor, 
        fill opacity=0.1, 
        line width=1.5pt
    ] plot[smooth cycle, tension=0.7] coordinates {
        (0, 0) (0.5, 2.5) (3.5, 2) (4.5, 0) (3, -2.5) (1, -1)
    };
    
    % Label for Synthetic
    \node[text=synColor, font=\bfseries] at (5.4, 2.2) {Synthetic Support ($\mathcal{S}_g$)};
    
    % ==========================================================
    % 3. DATA POINTS
    % ==========================================================
    
    % --- Real Points ---
    \foreach \i in {1,...,35} {
        \fill[realColor!60] 
        ({rand*1.5 - 0.5}, {rand*1.5}) circle (1.5pt);
    }
    
    % --- Good Synthetic Points (Overlap) ---
    \foreach \point in {(1, 1), (1.5, 0.5), (2, -0.5), (0.5, 0.5), (1.2, 1.5), (2.2, 0)} {
        \fill[synColor] \point circle (2pt);
    }
    \node[circle, draw=synColor, inner sep=2pt, thick] (goodpoint) at (1.5, 0.5) {};
    
    % --- Bad Synthetic Points (Hallucinations) ---
    \foreach \point in {(3.5, 1), (3.8, 0), (3.2, -1.5), (4, -0.5), (3.5, 0.5)} {
        \fill[synColor] \point circle (2pt);
    }
    \node[circle, draw=synColor, inner sep=2pt, thick] (badpoint) at (3.8, 0) {};
    
    % ==========================================================
    % 4. WASSERSTEIN VISUALIZATION (New Addition)
    % ==========================================================
    % Draw "Transport" arrows from Bad Points to Real Manifold
    % This implies "Moving" the probability mass
    \draw[-{Stealth}, transportColor, line width=1pt, dashed, shorten >=2pt, shorten <=2pt] 
        (3.5, 1) -- (2, 1.5);
    \draw[-{Stealth}, transportColor, line width=1pt, dashed, shorten >=2pt, shorten <=2pt] 
        (3.8, 0) -- (2.5, 0);
    \draw[-{Stealth}, transportColor, line width=1pt, dashed, shorten >=2pt, shorten <=2pt] 
        (3.2, -1.5) -- (1.8, -1.5);
    
    % Label for the arrows
    \node[text=transportColor, font=\footnotesize\bfseries, rotate=15] at (2.93, 0.63) {Transport Cost($W_1$)};

    % ==========================================================
    % 5. ANNOTATIONS
    % ==========================================================
    
    % Precision/Hallucination Labels
    \node[align=left, font=\footnotesize, text=synColor!80!black] (goodlabel) at (1, -4) 
    {\textbf{Realistic Sample}\\(In Support)};
    \draw[->, synColor, dashed] (goodlabel) -- (goodpoint);
    
    \node[align=left, font=\footnotesize, text=synColor!80!black] (badlabel) at (5, -3) 
    {\textbf{Hallucination}\\(Out of Support)};
    \draw[->, synColor, dashed] (badlabel) -- (badpoint);
    
    % Precision Box (Top Right)
    \node[align=center, fill=white, draw=gray, rounded corners, drop shadow, inner sep=5pt] at (5.2, 3.5) {
        
    };

    % Wasserstein Box (Bottom Left - NEW)
    \node[align=center, fill=white, draw=transportColor, rounded corners, drop shadow, inner sep=5pt] at (-3.5, -3.5) {
        \textbf{Wasserstein Distance}\\
        \textcolor{transportColor}{Cost to move $\mathcal{S}_g \rightarrow \mathcal{S}_r$}\\
        \scriptsize $W_1 = \inf \mathbb{E}[\|x - y\|]$
    };
    
    % Legend
    \matrix [draw, fill=white, below left, row sep=2mm, column sep=2mm] at (7.5, -3.5) {
        \fill[realColor!60] (0,0) circle (2pt); & \node[font=\scriptsize]{Real Data}; \\
        \fill[synColor] (0,0) circle (2pt); & \node[font=\scriptsize]{Synthetic Data}; \\
        \draw[-{Stealth}, transportColor, dashed, thick] (0,0) -- (0.5,0); & \node[font=\scriptsize]{Transport ($W_1$)}; \\
    };
    
\end{tikzpicture}
\caption{Manifold Validation (Wasserstein on Features)}
\end{figure}

\subsection{Downstream Performance (TSTR)}
\label{sec:tstr_results}

The central experiment trains the Apollo pipeline exclusively on CVAE-generated synthetic spectra and evaluates on held-out real patients. Results are presented in Table~\ref{tab:tstr}.

\begin{table}[H]
\caption{TSTR vs.\ TRTR Comparison (6-class glioma grading)}
\label{tab:tstr}
\begin{tabular*}{\textwidth}{@{\extracolsep\fill}lccc}
\toprule
Fold & TRTR Accuracy & TSTR Accuracy & $\Delta$ \\
\midrule
Fold 0 & 0.373 & --- & --- \\
Fold 1 & 0.263 & --- & --- \\
Fold 2 & 0.319 & --- & --- \\
Fold 3 & 0.288 & --- & --- \\
Fold 4 & 0.233 & --- & --- \\
\midrule
\textbf{Mean} & $\mathbf{0.295}$ & --- & --- \\
\bottomrule
\end{tabular*}
\note{TSTR columns to be filled after completing the synthetic generation and evaluation pipeline within each CV fold. Success criterion: TSTR accuracy within 5\% of TRTR.}
\end{table}

% TODO: Fill TSTR results after running notebook 05

% Latent Space Visualization:
% TODO: PCA plot of the 32-dimensional latent space for each class 

\subsection{Reconstruction Quality}
\label{sec:reconstruction}
The CVAE faithfully reproduces the major spectral peaks (phenylalanine at $1003$ cm$^{-1}$, CH$_2$ at $1450$ cm$^{-1}$) while smoothing high-frequency noise, consistent with the information bottleneck hypothesis.

% TODO: Include reconstruction figure generated by notebook 05
% \begin{figure}[H]
%     \includegraphics[width=\textwidth]{figures/reconstruction_comparison.pdf}
%     \caption{Spectral reconstruction for representative samples. Top: real (black). Bottom: CVAE reconstruction (red, mirrored).}
%     \label{fig:recon}
% \end{figure}

\subsection{Privacy and Diversity Metrics}
\label{sec:privacy}
Following \cite{alaa2022faithful}, we evaluate:
\begin{itemize}
    \item \textbf{Authenticity:} Minimum Euclidean distance from each synthetic sample to its nearest real neighbor. High distances indicate the model generates novel spectra rather than memorizing training examples.
    \item \textbf{$\beta$-Recall:} Coverage of the real data manifold by synthetic samples. Low recall indicates mode collapse.
    \item \textbf{$\alpha$-Precision:} Fraction of synthetic samples within the support of the real distribution. Low precision indicates hallucinated spectra.
\end{itemize}
% I highly suspect that those metrics between synthetic and real spectra will yield poor results given the high-dimensionality and noise in the real data and the denoising effect of the VAE, but they are still important to report for completeness and to contextualize the TSTR results. If TSTR accuracy is high despite low precision/recall, it would suggest that the CVAE is learning a simplified, denoised manifold that captures the essential diagnostic features while ignoring patient-specific noise.
% TODO: Fill with actual metric values from notebook 05

\section{Discussion}
\label{sec:discussion}

\subsection{The Generalization Gap}
\label{sec:gap_discussion}
The TRTR baseline accuracy of $\sim$30\% for the 6-class task---despite being above chance (16.7\%)---quantifies the severity of the domain shift problem in Raman spectroscopy. This result is consistent with the hypothesis that batch effects (spectrometer calibration, tissue preparation, laser focus depth) contribute more variance than the subtle biochemical differences between glioma subtypes. Our results are consistent with prior work, reporting roughly 80\% accuracy for binary classification (tumor vs. healthy or mutant vs. wild-type).

Our strict patient-level cross-validation reveals the true clinical challenge: building a model that generalizes to \emph{unseen patients} from \emph{unseen acquisition sessions}.

\subsection{The Role of Spatio-Spectral Annotation}
The unsupervised GMM-based annotation step is critical for two reasons. First, it removes the confounding signal of healthy tissue, which would otherwise dominate the learned latent space. Second, by operating per-slide with spatial context, it avoids the failure mode of global clustering where batch effects would drive the cluster assignments rather than biological differences.

\subsection{Interpretation of TSTR Results}
The ``Train on Synthetic, Test on Real'' (TSTR) evaluation tests whether the CVAE has learned the true underlying biochemical distribution of glioma subtypes. If a classifier trained \textit{solely} on synthetic data can diagnose real patients with accuracy comparable to the TRTR baseline, it demonstrates that the synthetic spectra contain the necessary diagnostic features (e.g., phenylalanine, amide I/III bands) without relying on patient-specific artifacts.

Critically, even if the TSTR accuracy only \emph{matches} the TRTR baseline ($\sim$30\%), this would still be a significant finding: it would prove that the generative model learned the complex 6-class biology while discarding patient-specific noise.

\subsection{Privacy and Data Sharing}
A significant advantage of this framework is privacy preservation \cite{kaissis2020secure}. Raman spectra, like other high-dimensional biomedical signals, can theoretically serve as biometric identifiers. Synthetic data with sufficiently large distance from the nearest real neighbor breaks this one-to-one link, enabling multi-institutional data sharing without navigating complex legal requirements under GDPR and HIPAA.

\subsection{Limitations and Future Work}
Several limitations should be noted:
\begin{itemize}
    \item \textbf{Gaussian Prior:} The standard normal prior assumption in VAEs can lead to overly smooth reconstructions, potentially attenuating high-frequency, low-intensity peaks with diagnostic value. Diffusion Probabilistic Models may better capture these fine-grained details.
    \item \textbf{Class Imbalance:} The dataset exhibits significant class imbalance across the 6 LGm subtypes. While the CVAE can generate balanced synthetic datasets by construction, the encoder's classification head may still be biased toward majority classes.
    \item \textbf{Single Dataset:} All experiments are conducted on a single institutional dataset. Multi-site validation would strengthen the claims about domain shift robustness.
    \item \textbf{Architectural Exploration:} Moving the auxiliary classification head post-latent space (creating a supervised information bottleneck) may further decouple convolutional bypass paths from the latent representation.
\end{itemize}
\section{Conclusion}
\label{sec:conclusion}

This study addresses the critical challenge of data scarcity in Raman-based glioma diagnostics through a three-stage framework: (1) unsupervised spatio-spectral annotation to isolate tumor tissue from mixed biopsy scans, (2) a $\beta$-Conditional Variational Autoencoder to learn class-conditional spectral distributions in a compressed latent space, and (3) a rigorous TSTR evaluation protocol with strict patient-level data isolation.

Our key findings are:
\begin{enumerate}
    \item The TRTR baseline for 6-class glioma grading under strict patient-level cross-validation is approximately 30\%, quantifying the severity of the generalization gap caused by batch effects.
    \item The per-slide GMM annotation pipeline successfully isolates $\sim$88,000 pure tumor spectra from $\sim$270,000 total tissue spectra without requiring manual labeling.
    \item The CVAE architecture with auxiliary classification produces faithful spectral reconstructions while organizing the latent space by tumor subtype.
\end{enumerate}

These results establish the methodological foundation for synthetic data augmentation in spectral diagnostics. If TSTR accuracy approaches the TRTR ceiling, it would demonstrate that generative models can learn the true biochemical signatures of glioma subtypes while filtering out the patient-specific noise that currently prevents clinical deployment of AI-assisted intraoperative Raman spectroscopy.

\bibliographystyle{plain}
\bibliography{references}

\end{document}